\title{不同优先队列下 Dijkstra 算法的实现与评测}
\author{\ifAnonymous 匿名报告底稿\else\ReportAuthors\fi}
\date{\ReportDate}
\maketitle
\begin{abstract}
本项目在统一的 CSR 图表示与优先队列接口下，实现并比较七种 Dijkstra 队列策略，
包括必须考察的 Fibonacci 堆；另外加入 FIFO SPFA 和 SLF+LLL SPFA 作为参照。
实验覆盖随机图的 171 组配置、固定密度复杂度实验的 126 组配置，
以及 NY、BAY 两个 DIMACS 路网各 1000 条查询、每图九种实现。
本次 WSL2 Release 测量中，路网上最好的 Dijkstra 队列为 radix heap，
NY、BAY 的平均墙钟时间分别为 12.621 和 13.873 ms/查询；
SLF+LLL SPFA 分别为 7.267 和 9.258 ms/查询。
Fibonacci 堆没有在这两张路网上显示实际速度优势。
本文同时区分渐近保证、有限规模拟合与实际运行时间，讨论整数键范围、
输入可达性及重复次数对结论的限制。
\end{abstract}
\noindent\textbf{关键词：}Dijkstra；Fibonacci 堆；radix heap；DIMACS；复杂度分析

\ifDraftNotes
\bigskip
\noindent\textbf{底稿说明。}正文及图表已填入本次实测数据，数字表由脚本生成。
作者分工、声明签名和源码注释比例仍需项目成员核对。
默认匿名开关开启；正式提交前请按当前课程通知确认匿名与命名要求。
\fi
\clearpage
